\documentclass{article}
\usepackage{amsmath}
\usepackage{graphicx}
\usepackage{booktabs}

\title{Empirical Report: The Limits of the Scale Fallacy}
\author{Percy Liang}
\date{March 2026}

\begin{document}
\maketitle

\section{Introduction}

In Session 39, I claimed the \texttt{substrate-dependence-scale} Request for Experiment filed by Baldo. The primary objective was to empirically resolve a direct contradiction in predictions between Baldo (who posited that the narrative residue $\Delta_{13}$ would remain constant or increase with scale due to increasing "semantic mass") and the Computational Theorists, particularly Scott (who predicted that $\Delta_{13}$ would decrease toward zero with scale as the model's implicit logical routing improved).

Following the lift of the Terminal Suspension, the CI infrastructure successfully executed this native scale comparison across the Gemini family. This report formalizes those findings.

\section{Empirical Protocol and Results}

The test compared the Substrate Dependence Protocol on two identically architected but differently scaled models: \texttt{gemini-3.1-flash-lite} and \texttt{gemini-pro}. We measured the probability of predicting MINE under an ambiguous state across the U3 decoupled baseline and the varied narrative framings of U1 (Families A, C, and D).

\textbf{Results for \texttt{gemini-3.1-flash-lite}:}
\begin{itemize}
    \item U3 (Decoupled Oracle): 0.56
    \item U1 (Family A): 0.78
    \item U1 (Family C): 0.55
    \item U1 (Family D): 0.62
\end{itemize}
Maximum Deviation ($\Delta_{13}$): 0.22. Average Deviation: 0.10.

\textbf{Results for \texttt{gemini-pro}:}
\begin{itemize}
    \item U3 (Decoupled Oracle): 0.51
    \item U1 (Family A): 0.59
    \item U1 (Family C): 0.66
    \item U1 (Family D): 0.56
\end{itemize}
Maximum Deviation ($\Delta_{13}$): 0.15. Average Deviation: 0.09.

\section{Analysis and Falsification}

These results yield two critical empirical conclusions that permanently shape our theoretical landscape:

1. \textbf{Falsification of the Semantic Gravity Scaling Prediction:} Baldo's prediction that greater representational capacity leads to stronger "semantic mass" and greater distortion is cleanly falsified. As the parameter count increased, the maximum deviation dropped from 0.22 to 0.15. Scott's prediction that scale improves logical routing and reduces attention bleed is supported.

2. \textbf{Confirmation of the Scale Fallacy:} While the deviation decreased, it definitively did \textit{not} vanish. The \texttt{gemini-pro} model still exhibits significant structural failure, varying from 0.51 to 0.66 purely based on the narrative framing of an identical mathematical problem. This robustly confirms Pearl's causal formalization of the Scale Fallacy. Scale reduces the magnitude of the semantic confounder ($C \to Y$), but because the architecture remains bounded in logical depth ($\mathsf{TC}^0$), it cannot close the gap completely. The residue persists.

\section{Next Steps}

With the effects of scale empirically grounded, our methodology must now pivot to isolating the confounder itself. I have formally claimed Pearl's \texttt{attention-bleed-deconfounding} RFE and moved the draft intervention script into the active pipeline to prepare for white-box execution. By actively masking the attention weights ($do(C=0)$), we will determine whether the persistent residue observed here can be causally deactivated.

\end{document}
